\paragraph{Main relationships}
\mbox{}\par

An account can send or receive many friend requests, appear in many friendships
and blocks, and take part in many conversations. Where a relationship has two
roles, two attributes reference \texttt{UserAccount.id} with different meanings,
for example the sender and the receiver of a friend request.

\begin{itemize}
  \item \textbf{Accounts and friendships:}
    \texttt{FriendRequest.senderId} and \texttt{receiverId},
    the two attributes of the \texttt{Friendship} primary key,
    and \texttt{UserBlock.blockerId} and \texttt{blockedUserId}
    all reference the matching account.
  \item \textbf{Accounts and conversations:}
    \texttt{Conversation.creatorId} references the creator.
    The many-to-many relationship between accounts and conversations is
    resolved through \texttt{ConversationMember}, whose composite primary key is
    \texttt{(conversationId, userId)}.
  \item \textbf{Conversations and messages:}
    every message belongs to one conversation. The pair
    \texttt{(conversationId, senderId)} in \texttt{Message}
    references the primary key of \texttt{ConversationMember}, which ensures
    the sender has a membership record in that conversation.
    Whether the sender is still a member when sending is checked by the backend.
  \item \textbf{Messages and personal history:}
    \texttt{HiddenMessage} links a user to a message.
    The composite primary key \texttt{(userId, messageId)} prevents recording
    the same hidden message twice for that user.
  \item \textbf{Messages and spam reports:}
    \texttt{SpamReport} references the reporter, the reported account
    and the related message; a message can appear in several reports.
  \item \textbf{Accounts and activity:}
    each \texttt{LoginHistory} or \texttt{AppOpenEvent} record
    belongs to one account; an account can have many records over time.
  \item \textbf{Encryption keys:}
    each account has at most one \texttt{UserKey}. A group has many
    \texttt{GroupKey} versions, keyed by \texttt{(conversationId, keyVersion)};
    each version has one \texttt{GroupKeyShare} per member who may read it.
    \texttt{Message.(conversationId, keyVersion)} references the key version
    used to encrypt the message.
  \item \textbf{Group events:}
    a group has many \texttt{GroupEvent} records; \texttt{actorId} and
    \texttt{targetUserId} reference the accounts involved.
  \item \textbf{Account locks:}
    \texttt{AccountLockLog} references the affected account, the acting
    administrator and, optionally, the \texttt{SpamReport} that led to the lock.
\end{itemize}

\paragraph{Planned constraints and data rules}
\mbox{}\par

Primary keys, foreign keys, unique constraints and \texttt{CHECK} constraints are
implemented in \texttt{script/create\_database.sql}; for example, a message's sender
must be a member of its conversation, a key share can only go to a member, and a
spam report must point to a message written by the reported account.
Rules about access rights, membership status or changing several records at once
are also checked by the backend; the keys in the ERD alone are not enough for
the whole business logic.

\begin{enumerate}
  \item \textbf{Accounts:} the username and email are unique;
    passwords are stored only as hashes. \texttt{systemRole} separates
    system Admins from Users; \texttt{status} separates active accounts
    from locked ones.
  \item \textbf{Friends and blocks:} the two people in a relationship must differ.
    \texttt{Friendship} stores the ID pair with the convention
    \texttt{userId1 < userId2}, combined with the composite primary key so both
    directions of the same relationship are never stored. A block is directed:
    A blocking B does not mean B blocks A.
  \item \textbf{Conversations and members:} \texttt{type} is
    \texttt{PRIVATE} or \texttt{GROUP}; the group name is used for group conversations.
    A private conversation has two participants; creating a group needs the creator
    and at least one other user. The composite primary key of a member prevents
    the same person from appearing twice in a conversation.
  \item \textbf{Group rights:} \texttt{memberRole} is
    \texttt{ADMIN} or \texttt{MEMBER}, independent of system administrator rights.
    Only group admins can remove members, as the assignment requires. On removal the
    backend sets \texttt{leftAt} instead of deleting the \texttt{ConversationMember}
    row, so the composite foreign key of \texttt{Message} stays valid and the person's
    earlier messages remain; when that person is added again, the same row is reused
    and \texttt{leftAt} is cleared.
    When a message is sent, the backend checks the current membership and rights.
  \item \textbf{Locked accounts:} when \texttt{UserAccount.status} is \texttt{LOCKED},
    nobody can message the account privately, send it or accept its friend requests,
    block it, add it to a group or make it a group admin; others can only unfriend it,
    decline its friend requests and, as group admins, remove it from a group.
  \item \textbf{Message status:} the \textbf{Sent}, \textbf{Waiting} and
    \textbf{Delivered} status of a private message is derived from
    \texttt{Message.deliveredAt}, with no extra status column; group messages only
    show \textbf{Sent}.
  \item \textbf{Spam reports:} the reporter must be allowed to see the related
    content; the reported message must belong to the reported account.
    The backend checks these conditions on top of the foreign keys.
  \item \textbf{Encryption and logs:} a key share can only go to someone with a
    membership record in the group, and a message can only name a key version of
    its own conversation (foreign keys); the backend also ensures that keys exist
    only for group conversations, that \texttt{Message.keyVersion} stays empty in
    private conversations and that a new key version is shared only with current
    members. \texttt{AccountLockLog.action} is \texttt{LOCK} or \texttt{UNLOCK},
    \texttt{GroupEvent.type} takes only the listed values, and a user can report
    the same message only once (check and unique constraints).
  \item \textbf{Consistency:} accepting a friend request updates the request and
    creates the friendship in the same transaction. Blocking and unfriending also
    apply their related changes together to avoid contradictory states.
\end{enumerate}

\paragraph{Clearing history and group encryption}
\mbox{}\par

The proposed design treats clearing history as removing it for the person who does
it. \texttt{HiddenMessage} records each hidden message;
\texttt{ConversationMember.historyClearedAt} marks the point before which the whole
history is cleared for that member. When messages are read, hidden messages and
messages before that point are not shown to that person; other participants keep
their history. The team confirms this interpretation when the design is finalised.

Group encryption follows the end-to-end model of encrypted group chats in
Facebook Messenger~\cite{ref:messenger-e2ee}: the server stores and forwards only ciphertext and never
holds a group key in plain form.

\begin{enumerate}
  \item \textbf{User keys:} at registration the client creates an RSA key pair~\cite{ref:rfc8017}.
    The public key is stored in \texttt{UserKey.publicKey}; the private key is
    encrypted on the client with a key derived from the password (PBKDF2~\cite{ref:rfc8018} with
    \texttt{keySalt}), so the user can sign in from any computer.
  \item \textbf{Turning encryption on:} the client of a group admin creates a
    random AES-256 group key (version 1), wraps it with every member's public key
    and uploads one \texttt{GroupKeyShare} per member;
    \texttt{Conversation.isEncrypted} becomes true.
  \item \textbf{Sending and reading:} the sender encrypts the text with the latest
    key version (AES-GCM~\cite{ref:nist-gcm}, a new IV per message) and stores the version in
    \texttt{Message.keyVersion}; members unwrap the matching share with their
    private key to decrypt.
  \item \textbf{Membership changes:} a new member receives a share of the current
    version. When a member leaves or is removed, an admin's client creates a new
    version and shares it only with the remaining members, so the former member
    cannot read later messages.
  \item \textbf{Turning encryption off:} new messages are stored as plain text with
    an empty \texttt{keyVersion}; earlier messages stay encrypted.
\end{enumerate}

Because the server cannot read encrypted messages, searching them is done on the
client after decryption, the AI assistant is not available in encrypted groups,
and a spam report on an encrypted message carries a plain-text copy in
\texttt{SpamReport.reportedContent} sent by the reporter. Changing a password only
re-encrypts the private key. A password reset (a random password sent by email)
cannot recover the old private key: the user receives a new key pair and can read
older encrypted history again only after a group admin shares the keys with the
new public key.

Deleting an account is a soft delete: the administrator's action sets
\texttt{status} to \texttt{DELETED} and records \texttt{deletedAt}, and no row is
removed. Messages, group history, spam reports and lock history therefore stay
complete, every foreign key to \texttt{UserAccount} rejects physical deletion, and a
deleted account can no longer sign in or be found. Its username and email may be
registered again, because uniqueness only applies to accounts that are not deleted.

\paragraph{Data for statistics}
\mbox{}\par

The model has no separate statistics tables when the figures can be computed from
the source data. The Admin functions map as follows:

\begin{itemize}
  \item \hyperlink{F14}{F14}: reads sign-ins from
    \texttt{LoginHistory}, joins the account profile and sorts by time.
  \item \hyperlink{F17}{F17}--\hyperlink{F18}{F18}: filters
    \texttt{UserAccount.createdAt} by the selected year; counts new accounts
    per month of that year.
  \item \hyperlink{F19}{F19}: counts relationships in \texttt{Friendship};
    \textbf{Total network} is the union of direct friends and friends of friends,
    excluding the account itself and counting each account once.
    \textbf{Friends of friends only} also excludes direct friends; a mutual friend
    is found through two consecutive friendships.
  \item \hyperlink{F20}{F20}: counts app opens from
    \texttt{AppOpenEvent}; checks messages, conversations and members to find
    how many people and groups were chatted with in the selected year.
  \item \hyperlink{F21}{F21}: counts distinct accounts that opened the app in
    each month, not the total number of opens.
\end{itemize}

Sign-in history and app opens are kept apart because reopening the app is not
necessarily a new sign-in. Online/offline status will be determined by the backend
from live connections; an old sign-in record does not prove a user is online.

The \textbf{Period} picker applies one calendar year to tables and charts.
A yearly query covers the start of the year up to the start of the next year in the
system's time zone. Charts always show 12 months; the yearly total of
\textbf{Active users} counts each account once, even when it was active in several
months.

\paragraph{Data for administration extensions}
\mbox{}\par
\texttt{UserAccount.status} holds the current \texttt{ACTIVE}/\texttt{LOCKED}
state, and \texttt{AccountLockLog} keeps its history: each lock or unlock records
the account, the administrator, the time and, when it comes from the spam report
screen, the report. The \textbf{Prior locks} count and the \textbf{Admin lock
history} on the Spam Reports screen read from this table. The backend checks that
\texttt{adminId} belongs to an account with \texttt{systemRole = ADMIN} and
updates the status and the log in one transaction.

The \textbf{Activity} and \textbf{Monitoring} tabs of Chat Groups combine two
sources. Message rows come from \texttt{Message} (an encrypted message is shown
only as ``sent an encrypted message''). Every other change is a
\texttt{GroupEvent} written by the backend in the same transaction as the change
itself: creating the group, renaming it, adding a member, a member leaving or
being removed, a role change, and turning encryption on or off.
\texttt{ConversationMember} alone cannot provide this, because it keeps one row per
member, cannot tell leaving from being removed, and overwrites \texttt{joinedAt}
when a member is added again.
